\section{Admission orders and confluence}
\label{sec:admission}

This section proves Theorem~B. A theory is often extended by admitting items one at a
time: definitions, axioms, components, each available once none of its prerequisites is still
pending.
Whether the result depends on the order of admission is a confluence question, and the
classical tools answer it.

\subsection{Rewriting and Newman's lemma}

\begin{definition}[Rewriting notions]
\label{def:rewriting}
Let $\to$ be a relation on a set $S$ and $\to^{*}$ its reflexive and transitive closure.
Two states $y,z$ are \emph{joinable} if $y\to^{*}w$ and $z\to^{*}w$ for some~$w$. The
relation is \emph{locally confluent} if $x\to y$ and $x\to z$ imply that $y,z$ are
joinable, and \emph{confluent} if $x\to^{*}y$ and $x\to^{*}z$ imply that $y,z$ are
joinable. It \emph{terminates} if there is no infinite chain $x_0\to x_1\to\cdots$, that
is, if the converse relation is well founded. A state $x$ is \emph{normal} if $x\to y$
for no~$y$.
\end{definition}

\begin{theorem}[Newman's lemma]
\label{thm:newman}
A terminating, locally confluent relation is confluent.
\end{theorem}

\begin{proof}
By well-founded induction on $x$ along the converse of $\to$, we show that any two paths
$x\to^{*}y$ and $x\to^{*}z$ end in joinable states. If either path is empty the claim is
immediate. Otherwise the paths begin $x\to a\to^{*}y$ and $x\to b\to^{*}z$. By local
confluence $a\to^{*}w$ and $b\to^{*}w$ for some~$w$. The induction hypothesis at $a$
joins $y$ and $w$ at some $v$, and at $b$ joins $z$ and $w$ at some $t$. Applied once more
at $a$ to the paths $a\to^{*}w\to^{*}v$ and $a\to^{*}w\to^{*}t$, it joins $v$ and $t$ at
some $u$, and $u$ joins $y$ and~$z$.
\end{proof}

\begin{corollary}[Unique normal forms]
\label{cor:normal-forms}
If $\to$ is confluent, any two normal states reachable from a common state are equal. If
$\to$ also terminates, every state reaches exactly one normal state.
\end{corollary}

\begin{proof}
Two normal states reachable from one state are joinable, and a path out of a normal state
is empty, so they are equal. Under termination every state reaches some normal state, by
well-founded induction: either it is normal, or it steps to a state that reaches one.
\end{proof}

For a finite system, local confluence can be decided mechanically.

\begin{proposition}[Finite decision of local confluence]
\label{prop:finite-peaks}
Let $S$ be covered by a finite list, with decidable equality and a decidable one-step
relation. Then reachability $x\to^{*}y$ holds exactly when $y$ lies in every subset of the
list that contains $x$ and is closed under steps, and this can be checked by enumerating
the subsets. Joinability is decided by searching the list for a common reachable state,
and local confluence by checking every one-step peak.
\end{proposition}

\begin{proof}
The set of states reachable from $x$ is itself closed under steps and contains $x$, so a
state in every such subset is reachable. Conversely, a closed set containing $x$ contains
everything reachable from $x$, by induction on the path. The rest is a finite search over
the covering list.
\end{proof}
\begin{remark}[Efficient search]
\label{rem:graph-search}
The subset enumeration is exponential in the length of the list. It is the form the Lean
proof checks, because the characterization by closed subsets makes its correctness easy
to verify. In practice reachability is decided by graph search: a breadth-first or
depth-first search from $x$ visits each state at most once and computes the set of states
reachable from $x$ with at most $n^2$ tests of the step relation, for a list of length~$n$. Local
confluence is then checked by one such search from each end of each one-step peak.
\end{remark}

\begin{example}[Termination is needed]
\label{ex:four-state}
On four states $a,b,c,d$ let the only steps be $b\to a$, $b\to c$, $c\to b$ and
$c\to d$. Every one-step peak is joinable: the peak $b\to a$, $b\to c$ joins at $a$ via
$c\to b\to a$, and the peak $c\to b$, $c\to d$ joins at $d$ via $b\to c\to d$. The relation
is not confluent, since $b\to a$ and $b\to c\to d$ reach the distinct normal states $a$
and~$d$. The relation does not terminate: $b\to c\to b\to\cdots$.
\end{example}

\subsection{Admission systems}

\begin{definition}[Admission system]
\label{def:admission-system}
An \emph{admission system} over a set $I$ of items consists of a prerequisite relation
$p\prec i$ on $I$ and an audit map from finite lists of items to a set of labels. A state
is the list of items still pending. An item $i$ is \emph{enabled} in a state $s$ if $i$ is
pending and no prerequisite of $i$ is pending, and the admission step removes an enabled
item: $s\to s\ominus i$, where $s\ominus i$ deletes the first occurrence of $i$ from the
list~$s$.
\end{definition}

The model tests only that no prerequisite is still pending. A prerequisite that was never
in the initial list therefore counts as satisfied. When every prerequisite of every
pending item lies in the initial list, ``no prerequisite of $i$ is pending'' means that
every prerequisite of $i$ has been admitted. Removing an item never makes an item pending
again. This monotonicity, that admission can only
enable, is what the next theorem uses.

\begin{theorem}[Admission order does not matter]
\label{thm:admission}
In every admission system:
\begin{enumerate}[nosep]
  \item every admission run terminates;
  \item if $s\to y$ and $s\to z$, then $y=z$ or there is $w$ with $y\to w$ and $z\to w$;
  \item the admission relation is confluent;
  \item from any state $s$, every maximal admission order reaches the same terminal state,
        and so the same audit value.
\end{enumerate}
If the initial list has no repetitions, neither does any state reached from it.
\end{theorem}

\begin{proof}
Termination and one-step commutation give confluence by Newman's lemma, and uniqueness of the terminal state
follows.
(1) Each step shortens the pending list. (2) Let $y=s\ominus i$ and
$z=s\ominus j$ with $i\neq j$. Removing $j$ keeps $i$ pending and pending no new
item, so $i$ is still enabled in $z$; likewise $j$ is enabled in $y$. Both orders then
reach $(s\ominus i)\ominus j=(s\ominus j)\ominus i$. (3) By (2) the relation is locally confluent, and by (1) and
\Cref{thm:newman} it is confluent. (4) This is \Cref{cor:normal-forms}; equal terminal
states have equal audit values. The last statement holds because removing an element
preserves the absence of repetitions.
\end{proof}

The terminal state need not be empty. Items on a cycle of prerequisites, or depending on
such items, are never enabled and remain pending. The theorem says they remain pending in
every order.

\begin{example}[One disabling guard]
\label{ex:disabling}
Take items $a$ and $b$ and a guarded step in which $a$ may always be admitted but $b$ may
be admitted only while $a$ is still pending. From $\{a,b\}$ the relation terminates.
Admitting $a$ first reaches $\{b\}$, which is terminal because $b$ is now blocked.
Admitting $b$ first reaches $\{a\}$ and then $\varnothing$. The two terminal states differ,
and so do their lengths. The relation is not locally confluent, since the peak
$\{a,b\}\to\{b\}$, $\{a,b\}\to\{a\}$ cannot be joined.
\end{example}

\Cref{fig:admission} shows this example beside that of \Cref{ex:four-state}.

\begin{figure}[htbp]
\centering
\begin{tikzpicture}[font=\footnotesize,>=Stealth,
  state/.style={draw=gray,rounded corners=2pt,minimum width=1.25cm,minimum height=0.65cm,inner sep=2pt}]
  \node at (3.05,0.15) {(a) One disabling guard};
  \node[state] (ab) at (3.05,-1.0) {$\{a,b\}$};
  \node[state] (onlyb) at (1.2,-2.55) {$\{b\}$};
  \node[state] (onlya) at (4.9,-2.55) {$\{a\}$};
  \node[state] (empty) at (4.9,-4.05) {$\varnothing$};
  \draw[->] (ab) -- node[above left] {admit $a$} (onlyb);
  \draw[->] (ab) -- node[above right] {admit $b$} (onlya);
  \draw[->] (onlya) -- node[right] {admit $a$} (empty);
  \node[below=3pt of onlyb] {terminal};
  \node[below=3pt of empty] {terminal};
  \draw[gray] (6.5,0.55) -- (6.5,-4.7);
  \node at (9.8,0.15) {(b) Nonterminating relation};
  \node[state] (b) at (8.45,-1.7) {$b$};
  \node[state] (c) at (11.15,-1.7) {$c$};
  \node[state] (a) at (7.75,-3.55) {$a$};
  \node[state] (d) at (11.85,-3.55) {$d$};
  \draw[->] (b) to[bend left=20] (c);
  \draw[->] (c) to[bend left=20] (b);
  \draw[->] (b) -- (a);
  \draw[->] (c) -- (d);
  \node[below=3pt of a] {normal};
  \node[below=3pt of d] {normal};
\end{tikzpicture}
\caption{(a) Admitting $a$ from $\{a,b\}$ leaves terminal $\{b\}$, while admitting $b$ first leads to terminal $\varnothing$. (b) The steps $b\to a$, $b\to c$, $c\to b$, and $c\to d$, with normal states $a$ and $d$.}
\label{fig:admission}
\end{figure}

The guard violates the monotonicity of admission: admitting $a$ disables~$b$.

\begin{example}[Admission without constraints]
\label{ex:free-admission}
With two items and no prerequisites, both orders admit both items and reach the empty
terminal state.
\end{example}

The catalog records a finite admission system as a list of critical-pair records, each
marked legal or not, joinable or not and audit-equivalent or not, and states that a
critical-pair criterion is equivalent to confluence at the declared scope. Both sides of
that equivalence are predicates on the same list of records, so the statement is a
reformulation rather than a confluence theorem (\Cref{sec:sep-definitional}).
\Cref{thm:admission} is the corresponding semantic result, and \Cref{ex:disabling} proves
that different orders from one initial state can reach different terminal states once an
admission may disable another.
